\section{Almost bounded values of the Collatz map}

\begin{definition}[Collatz map]
\label{def:collatz-map}
\lean{CollatzMapAlmostBoundedValues.collatz}
\leanok
The Collatz map sends an even $n$ to $n/2$ and an odd $n$ to $3n+1$.
\end{definition}

\begin{definition}[Attaining a lower value]
\label{def:collatz-below}
\lean{CollatzMapAlmostBoundedValues.AttainsBelow}
\leanok
\uses{def:collatz-map}
$\mathrm{AttainsBelow}(f,n)$ means that some Collatz iterate of $n$ is strictly less than $f(n)$.
\end{definition}

\begin{definition}[Logarithmic density zero]
\label{def:collatz-log-density}
\lean{CollatzMapAlmostBoundedValues.logWeightSum,CollatzMapAlmostBoundedValues.LogDensityZero}
\leanok
The exceptional set is measured using logarithmic density.
\end{definition}

\begin{theorem}[Almost all Collatz orbits attain almost bounded values]
\label{thm:collatz-almost-bounded}
\lean{CollatzMapAlmostBoundedValues.MainTheorem}
\notready
\uses{def:collatz-below,def:collatz-log-density}
For every $f(n)\to\infty$, the starting values whose Collatz orbit never falls below $f(n)$ have logarithmic density zero.
\end{theorem}
